\section{Πρόβλημα 5}

Θεωρούμε το παίγνιο συμφόρησης του Σχήματος~\ref{fig.congestion-graph}. Έστω ότι
έχουμε 17 παίκτες και κάθε παίκτης έχει να διαλέξει ανάμεσα σε 3 διαδρομές
$(A, B, C)$, για να πάει από ένα αρχικό σημείο $s$ σε έναν προορισμό $t$. Οι
καθυστερήσεις είναι $d_A(x) = 5x$, $d_B(x) = 7.5x$, $d_C(x) = 10x$, αντίστοιχα
και όπου $x$ είναι ο αριθμός των παικτών που επιλέγουν την αντίστοιχη διαδρομή.

\begin{figure}[h]
    \centering
    \includegraphics[width=0.6\textwidth]{congestion-graph}
    \caption{Το δίκτυο με τις τρεις διαδρομές A, B, C από $s$ σε $t$.}
    \label{fig.congestion-graph}
\end{figure}

Καταρχάς, παρατηρούμε ότι το παίγνιο είναι συμμετρικό. Δηλαδή, μας νοιάζει να
βρούμε μόνο τον αριθμό των παικτών που επιλέγουν κάθε διαδρομή, π.χ. αν έχουμε
$8$ παίκτες να επιλέγουν την διαδρομή $Α$ αυτοί μπορεί να είναι οποιοιδήποτε
$8$. Οπότε θέλουμε βρούμε τις τιμές $f_A, f_B, f_C$ που εκφράζουν τον αριθμό
των παικτών που επιλέγουν κάθε διαδρομή. Θέλουμε για ανά δύο διαδρομές $i, j$
να ισχύει ότι $d_i(f_i) \leq d_j(f_j + 1)$. Αν είχαμε συνεχείς τιμές για το
φορτίο κάθε διαδρομής, τότε θα θέλαμε
$d_A(f_A) = d_B(f_B) = d_C(f_C) \Leftrightarrow 5f_A = 7.5f_B = 10f_C$ και με
$f_A + f_B + f_C = 17$ θα παίρναμε
$f_A = \frac{102}{13}, f_B = \frac{68}{13}, f_C = \frac{51}{13}$. Αρκεί να
ελέγξουμε τις πλησιέστερες ακέραιες τιμές για τις οποίες ισχύει
$\forall i, j \in \{A, B, C\}, i \neq j: d_i(f_i) \leq d_j(f_j + 1)$. Για
$(f_A, f_B, f_C) = (8, 5, 4)$ έχουμε $d_A(8) = 40, d_B(5) = 37.5, d_C(4) = 40$
το οποίο αποτέλει σημείο ισορροπίας. Οποιαδήποτε άλλη παραπλήσια, ακέραια
τιμή (π.χ. $(f_A, f_B, f_C) = (7, 6, 4)$) θα παραβιάσει τις συνθήκες
ισορροπίας.\footnote{Ουσιαστικά ελέγχουμε το potential function
$\Phi(f_A, f_B, f_C) = \sum_{i\in\{A,B,C\}} \sum_{j=1}^{f_i} d_i(j)$ πάνω στο
επίπεδο $f_A + f_B + f_C = 17$. Η συνεχής λύση μας δίνει το μοναδικό ελάχιστο
λόγω της κυρτότητας από το γεγονός ότι οι συναρτήσεις είναι γνησίως αύξουσες.
Βρίσκουμε το ελάχιστο για ακέραιες τιμές δοκιμάζοντας να κινηθούμε προς κάθε μία
από τις κατευθύνσεις των αξόνων.}
